% IV-C · Redacta: Farid · Revisa: Jorly · Semanas 12 y 13 · ~0,3 pág.
% Debe cubrir: Figs. 4 y 5, el veredicto de H4 y la separación de pares
% (separacion_pares.csv).
\subsection{¿Por qué algunas categorías necesitan más? (H4)}
\pendiente{H4 y análisis por par}

\begin{figure*}[t]
\centering
\figura{fig_tsne.png}{0.85\textwidth}
\caption{\pendiente{t-SNE de los vectores GAP de M2-P al 30\,\% y al 100\,\% del dibujo.}}
\label{fig:tsne}
\end{figure*}

\begin{figure}[t]
\centering
\figura{fig_h4.png}{\columnwidth}
\caption{\pendiente{H4: confundibilidad frente a la mediana de $k^*/K$.}}
\label{fig:h4}
\end{figure}

% Fig. 6 (H5) y Fig. 7 (trazo a trazo) se definen aquí, después de H4, para
% que la numeración siga el plan; en IV-B se citan con \ref{fig:h5}.
\begin{figure}[t]
\centering
\figura{fig_h5.png}{\columnwidth}
\caption{\pendiente{H5: distribución de $K-k^*$ (trazos de anticipación).}}
\label{fig:h5}
\end{figure}

% Fig. 7 · Prepara: Jorly. Predicción de M1, M2 y M2-P en cada trazo de
% cuatro dibujos de los pares; respaldo de la demo en la exposición.
\begin{figure*}[t]
\centering
\figura{fig_trazo_a_trazo.png}{0.85\textwidth}
\caption{\pendiente{Predicción de cada modelo trazo a trazo en dibujos de los pares.}}
\label{fig:trazo}
\end{figure*}
